\section{Open decisions}
\label{sec:open}

Each one changes the entity set or the metrics, so an invented answer here would quietly
invalidate later work. Owners are named so that a decision has somewhere to land. Two items have
been resolved by the study, and one has a proposed ruling that the partner has not yet agreed to;
the table marks both, and the reasoning follows it so that a decision and its grounds stay
traceable to one place.

\begin{spectab}{|>{\raggedright\arraybackslash\bfseries}p{0.065\textwidth}|L{0.315\textwidth}|L{0.16\textwidth}|L{0.30\textwidth}|}
\hline
\rowcolor{headcell}\textbf{ID} & \textbf{Decision} & \textbf{Owner} & \textbf{Blocks} \\ \hline
\endfirsthead
\hline
\rowcolor{headcell}\textbf{ID} & \textbf{Decision} & \textbf{Owner} & \textbf{Blocks} \\ \hline
\endhead
OD-1 & Who is the data controller on paper: the City Health Office, or the Department of
  Health regional office? & CHO & The agreement, and every claim about who may
  publish. \\ \hline
\rowcolor{rowshade}OD-2 & What is the reporting period: daily, weekly or monthly? &
  CHO / RRL & Demand labels, the surge window, and whether a few hundred records is enough
  history. \\ \hline
OD-3 & Is distribution in scope, meaning stock positions and allocation recommendations, or
  is the study classifying only? & resolved & Whether two domain entities become
  core, and whether a shortage outcome is required. \\ \hline
\rowcolor{rowshade}OD-4 & How is a surge defined: facility-relative percentile, or
  district-wide outbreak signal? & study (proposed) & The label behind every reported metric.
  Needs partner agreement, since an existing system is likely to already define it. \\ \hline
OD-5 & Where does the aggregator run: university-hosted or CHO-hosted? & CHO / study &
  The strength of the privacy claim and the operational burden. \\ \hline
\rowcolor{rowshade}OD-6 & Are national facility codes and PSGC barangay codes obtainable
  for the pilot sites? & CHO & Entity identity; without them, facility identity must be
  redesigned. \\ \hline
OD-7 & Which quality metric is primary: a scaled error measure on demand
  counts, or balanced accuracy on surge days? And what threshold makes a model usable? &
  resolved & Every reported result, and the meaning of $\tau$ (Section~\ref{sec:metrics}). \\ \hline
\rowcolor{rowshade}OD-8 & In deployment, how is $Q_r$ computed without moving rows, when only
  parameters may cross to the aggregator? & CHO / adviser & What $Q_r$ means outside simulation: a
  per-facility score is neither a row nor a parameter, so the parameters-only promise may need
  narrowing (Section~\ref{sec:deployment}). \\ \hline
\end{spectab}

\noindent
Until an item is resolved, the text that depends on it should cite this table rather than assume
a default. A fallback stated as conditional is not an assumption: it applies only if the item is
still open after the review, and the section that states it says so.

\subsection{Resolved}

\textbf{OD-3 --- distribution is out of scope (resolved 2026-11-20).}
The study classifies surge periods at each facility. It does not forecast quantities, track stock
positions, or recommend allocations. The domain entities for distribution remain defined so that
the answer can change without redesigning the record, but no shortage outcome is required and none
is claimed.

\textbf{OD-7, metric --- balanced accuracy is primary (resolved 2026-11-20).}
$Q_r$ is balanced accuracy on surge periods. A low-volume facility's series is zero-heavy, and a
percentage error is undefined or explosive on zeros, whereas balanced accuracy never divides by a
count that may be zero. The comparison is also per facility, and an error measure is not
scale-free: ranking a station holding a few hundred records against one holding a few thousand
would need a scaling argument before any ranking could be stated.

\textbf{OD-7, threshold --- $\tau$ is the local-only baseline (resolved 2026-11-20).}
$\tau$ is the balanced accuracy of a model trained on one facility's own records, computed once per
experimental condition from the frozen baselines before any policy comparison runs, and declared
with that condition. $T^{*}$ then answers something a facility can act on: how long until
federation catches up to what it could already do alone. A condition whose local-only baseline
falls at or below $0.5$ --- what an always-negative classifier scores --- has no usable $\tau$ and
is reported as unevaluable rather than compared.

\subsection{Proposed --- awaiting partner review}

\textbf{OD-4 --- surge is facility-relative (proposed 2026-11-20, not agreed).}
A surge is a facility's own demand rising above its own past, because barangays do not surge
together: Barangay A may need a commodity when Barangay B does not. The definition is ready and is
what the study implements unless it is overridden. It should be read as contested rather than
settled: a partner that already runs a system will tend to hold its existing definition as the
correct one, and a district-wide outbreak signal is a defensible alternative. The disagreement is
not cosmetic, because the label changes every reported metric.

Two parts of the definition separate, and only one of them is negotiable.

\textbf{Negotiable: what counts as elevated.} A facility-relative percentile, a district-wide
outbreak signal, or the partner's existing rule.

\textbf{Not negotiable: the threshold is computable from past data alone.} Whatever rule is chosen,
the label at period $t$ uses records up to $t-1$ and never a facility's whole history. Under the
proposed definition that is
\[
\text{surge}_t = 1 \quad\Longleftrightarrow\quad D_t > P_q\!\left(D_{t-w},\ldots,D_{t-1}\right)
\]
with window $w$ and percentile $q$ declared per condition. A district-wide rule is admissible only
if it can be written in the same form; otherwise the threshold would not be computable at the point
it is needed.

\noindent
The same rule in arithmetic, for a facility reporting monthly with $w = 12$ and $q = 80$. The
twelve counts recorded before this month, sorted into ascending order:

\[
38,\;40,\;41,\;43,\;45,\;46,\;47,\;49,\;52,\;55,\;58,\;60
\]

\noindent
The eightieth percentile sits between the ninth value and the tenth, eight tenths of the way along:

\[
P_{80} \;=\; 52 + 0.8 \times (55 - 52) \;=\; 54.4
\]

\noindent
so the label for a month is decided by a single comparison against $54.4$:

\[
D_t = 58 > 54.4 \;\Longrightarrow\; \text{surge}_t = 1,
\qquad
D_t = 47 \le 54.4 \;\Longrightarrow\; \text{surge}_t = 0 .
\]

\noindent
That example exposes something the symbolic form hides. $P_{80}$ of twelve values is not one of the
recorded counts, so the percentile needs a convention. Linear interpolation between the two
neighbouring order statistics is the convention used here; taking the tenth value outright would
set the threshold at $55$ and flip a month of exactly $55$ units. The threshold is a computed
number, not a count anyone wrote down, and the convention is part of the definition.

\noindent\textbf{In plain terms.} Each facility is measured against its own recent past. If this
month's demand is above what four months in five of its own recent months looked like, that month is
a surge for that barangay. Another barangay's numbers never enter the comparison, and neither does
any fixed number of units, so a station that is always busy and one that is never busy are each
judged against themselves.

Three consequences are stated rather than discovered later. $q$ sets the class balance instead of
measuring Iligan's: at $q = 80$ roughly one period in five is a surge at every facility by
construction. A facility-relative label is not the same as an operational shortfall --- a facility
can sit at its own ninetieth percentile and still be nowhere near a stockout --- which OD-3 puts out
of scope and which this study therefore does not measure. And the estimate is unstable at the period
length currently reported: informal talks put local reporting at monthly or quarterly, so a
percentile over a couple of dozen periods rests on very little history, and neither $w$ nor $q$ can
be declared until OD-2 fixes the period.
